\section{Reproducing IRASim}
\label{sec:reproduce}

\paragraph{Pipeline.}
IRASim~\cite{zhu2024irasim} predicts 16 frames from the first frame and 15 end-effector actions, each the motion between two frames in the end-effector frame of the earlier one, plus the gripper state.
The SDXL~\cite{podell2023sdxl} VAE encodes each $256\times320$ frame into a $4\times32\times40$ latent, and a diffusion transformer~\cite{peebles2023dit} denoises the latents of frames 1 to 15 with the clean latent of frame 0 as the condition; it is trained with the noise-prediction loss~\cite{ho2020ddpm} and sampled with 50 PNDM~\cite{liu2022pndm} steps.
Actions enter through adaptive layer norm: Frame-Ada modulates each frame by its own action and Video-Ada every frame by one embedding of the sequence, while the baselines LVDM~\cite{he2022lvdm} and VDM~\cite{ho2022vdm} condition on the whole sequence in latent and in pixel space.

\paragraph{Setup.}
We use the authors' code, checkpoints, and evaluation settings on 105 Bridge~\cite{walke2023bridgedata} test clips of 16 frames at 5\,Hz, the first episodes of each of 12 instruction verbs, and on 57 Language-Table~\cite{lynch2022interactive} clips.
Cosmos-Predict2.5 without actions, at 3\,min per clip on an A100, runs on a fixed subset of 24 of the Bridge clips, and all Cosmos videos are resampled to 16 frames at $256\times320$.
The initial noise of a clip depends on its batch, so seeds are samples, and Frame-Ada is the mean of two.

\paragraph{Differences from the paper.}
We do not compute FVD~\cite{unterthiner2018fvd}, which gives no per-clip score and ignores the commanded actions.
Our PSNR is lower than reported, with the same ranking of models (\cref{tab:bridge}).
The authors' code scores the same frames against the same real videos, and on 64 test clips drawn at random Frame-Ada reaches 24.85\,dB, within 0.43\,dB of the reported 25.28, so the first episodes of each verb that we score are harder than the average test clip.

% ---- review comments, kept for the record (text above rewritten for concision on 2026-11-02) ----
% (Turn 3) "We use the code, checkpoints, and evaluation settings of [20] at commit c72b6da, changing only the episode list, the seed, and the batch size": Don't you think this sentence is weird? In professional paper they do not write in too much details like this, I think we should keep this information in our work notes
% Claude: agreed; the commit and the changed arguments stay in the code README (scripts/setup/irasim.sh pins the commit), the paper keeps one sentence.
% (Turn 3) "24 core clips over 8 verbs, extended to 105 clips over 12 verbs": Don't you think we should use only the larger number, 105 clips, and if there are not-calculate-yet numbers on these 105 clips, just run a new A100 session and finish them?
% Claude: partly disagree. Every main comparison (Tabs. 3, 4, executability, agreement) already uses 105 clips. What is on 24 only: four corrupted-action sets, four extra seeds, VDM, Cosmos without actions, and the judges, about 900 extra videos (Cosmos without actions alone is 3 min per clip on an A100, about 30 CU for the rest of the 105) plus judge calls, so several A100 hours and a real risk of missing 18:00. A main set plus a fixed subset for costly ablations is common practice if it is stated, which the text now does. The perturbation set on 105 clips (420 IRASim videos, roughly 3 to 4 A100 hours, unmeasured) is the one extension worth it if you still want it; I asked in the chat.
% (Turn 3) "and we do not compute FVD [15], which needs far more videos than we generate": Is there any reason else that we do not use this metric? Or should we generate more videos to use it since now is just 10:41, the deadline is 18:00
% Claude: yes, two more reasons, now in the text: FVD gives one number per set, so no per-clip paired tests, and it ignores the commanded actions (a model that moves the arm the wrong way but looks real can get a good FVD), which is exactly our question. Generating more videos would not fix that, and our FVD on a few hundred videos would still not be comparable with the paper's on the full test set. So I recommend not computing it.
% (Turn 3) "rerunning with the same batches reproduces the videos exactly, but a different batch composition left only 1 of 24 clips unchanged, so we treat seeds as samples and report results per see": Do you mean our reported results run in different batch sizes? So we have to report the results in different seeds? I do not want to report resutls in differents seeds like this, if can I think we should average them just as we do in task lang_grasp, or if it is necessary, run mor eexperiments.
% Claude: not batch sizes: the batch size is the same; what changes is which clips share a batch (24 clips, then 105), and that changes the random noise each clip starts from, so a seed is just one random sample, not a fixed video. It does not change quality. Agreed on averaging: Tabs. 3 and 4 now show one Frame-Ada row, the mean of seeds 3407 and 3408 (std at most 0.12 dB PSNR on Bridge), as in lang_grasp. The other models have one seed each; the seed-to-seed spread is far smaller than the gaps between models, and paired tests now hold with either seed (one exception reported: Language-Table NDTW against LVDM, seed 3408, p = 0.09).
% (Turn 3) "Our PSNR is 0.9 to 1.5 dB below the reported values on Bridge and 0.9 to 1.8 dB below on Language-Table, with the same ranking of models (Tabs. 3 and 4)": If the information already in the table, then I think we do not need to redundantly report them in here, or maybe we should write in a more general style, with key information is "with the same ranking of models" or "the ranking of the models keeps unchanged"
% Claude: agreed; the numbers stay in the tables, the text keeps the ranking.
% (Turn 3) "the clip subset may explain the offset, which we have not verified": Oh my god if we have not verified yet so why don't you tell me ealier, just do it if this is an important information, else, just remove this "unsrure" line
% Claude: sorry, it was written as unverified in the task notes but I should have raised it with you. Now checked in the IRASim code at the commit we use (evaluate/compute_psnr_ssim.py, evaluate/util/generate_true_sample_video.py): their PSNR skips frame 0, averages frames 1 to 15, and compares with the raw dataset frames, not VAE reconstructions, exactly as ours. So the protocol is ruled out, and the text now says only what is verified. Proving that the clip choice is the cause would need generating Frame-Ada on a random sample of the full test set (Kaggle, about 2 h); asked in the chat.
